% Generálva: html/build_tex.js (o09.js). Kézzel ne szerkeszd, a forrás a html/data mappa.
\setcounter{chapter}{8}
\chapter{\texorpdfstring{GSM: kapacitás, hívásfogadás, SMS}{GSM: kapacitás, hívásfogadás, SMS}}\label{f9}
Klaszterek és időrések, és a roamingos kérdések szabályai. A roamingos válaszoknál minden elem mellett ott van, melyik szabályból jön.

\section{Alapok, egyszerűen}
\subsection{\texorpdfstring{A1. A hálózat elemei}{A1. A hálózat elemei}}\label{b9-elemek}
\par\noindent{\small\renewcommand{\arraystretch}{1.25}\begin{tabularx}{\linewidth}{@{}>{\raggedright\arraybackslash}X>{\raggedright\arraybackslash}X@{}}\toprule
\textbf{Elem} & \textbf{Szerep}\\\midrule
MS & mobilállomás (készülék + SIM)\\
BTS & bázisállomás: a rádiós adó-vevő egy cellában\\
BSC & bázisállomás-vezérlő: több BTS-t irányít, a hívásátadás egy részét intézi\\
MSC & mobil kapcsolóközpont: hívásirányítás, kapcsolás\\
GMSC & átjáró MSC: a hálózatba kívülről érkező hívást ez fogadja, és ez kérdezi a HLR-t\\
HLR & honi adatbázis: az előfizető adatai és hogy éppen melyik VLR területén van\\
VLR & látogatói adatbázis: a területén éppen tartózkodók ideiglenes adatai\\
AuC & hitelesítési központ: a SIM titkos kulcsával (Ki) azonosít\\
EIR & készülék-nyilvántartás: IMEI-listák (fehér, fekete, szürke)\\
SMSC & SMS-központ: tárolja és továbbítja az üzeneteket\\
\bottomrule\end{tabularx}}\par

Az $\alpha$ és $\beta$ két szolgáltató hálózata, az index mutatja, melyiké az elem.

\subsection{\texorpdfstring{A2. A szabályok}{A2. A szabályok}}\label{b9-gsm}
\begin{enumerate}
\item \textbf{Hívás:} a hívott szám a hívott honi hálózatához tartozik, ezért a hívás ennek GMSC-jéhez fut, és ez a \textbf{hívott honi HLR}-jét kérdezi. A HLR tudja, melyik VLR területén van a hívott, onnan kér barangolási számot.
\item \textbf{SMS:} a \textbf{küldő} honi SMSC-je fogadja, tárolja és kézbesíti.
\item \textbf{Készülékellenőrzés:} a \textbf{kiszolgáló} (ahol a készülék éppen van) hálózat MSC/VLR-je ellenőrzi az IMEI-t a saját EIR-jénél.
\item \textbf{Mozgás beszélgetés közben:} hívásátadás (handover), ebben csak a BTS, a BSC és az MSC vesz részt. Helyzetfrissítés (és a VLR, HLR értesítése) csak hívás nélküli mozgásnál van.
\item Egy esemény (hívásindítás, bontás) először azon az elemen jelenik meg, amelyikkel a kezdeményező készülék közvetlenül beszél: a \textbf{saját BTS-én}.
\end{enumerate}

Kapacitás: a GSM 900 sávjában 124 vivő van, vivőnként 8 időrés. Egy $N$ cellás klaszter az összes vivőt egyszer használja fel.

% ------------------------------------------------------------------
\section{\texorpdfstring{9.1. feladat: Budapest: hány klaszter, hányan beszélhetnek?}{9.1. feladat: Budapest: hány klaszter, hányan beszélhetnek?}}\label{o9_01}
\setcounter{lepes}{0}
\begin{feladat}{9.1. feladat}
Budapest 525 km\textsuperscript{2}, FRF = 12, 525 cella. a) Hány klaszter? b) Hányan beszélhetnek egyszerre? c) Egy átlagos cellában hányan egy szolgáltatónál (kb. 40 csatorna)?

{\footnotesize\sffamily Forrás: 5-ösöm 25. oldal $\cdot$ 5o\_25\_1. Címkék: vizsga}
\end{feladat}
\begin{keresunk}
\begin{itemize}
\item a) a klaszterek száma
\item b) az egyszerre lehetséges beszélgetések száma
\item c) egy cella egy szolgáltatójánál egyszerre beszélők száma
\end{itemize}
\end{keresunk}
\begin{jeloles}
\par\smallskip\textbf{Adott még:}
\begin{itemize}
\item 525 cella (kb. 1 km\textsuperscript{2}-enként egy)
\item 124 vivő, vivőnként 8 időrés
\item c) egy szolgáltató kb. 40 vivőt használ
\end{itemize}
\end{jeloles}

\lepes{a) Klaszterek}
Minden klaszter 12 cellából áll.

\begin{center}\adjustbox{max width=\linewidth}{$\displaystyle \frac{525}{12}=43{,}75\approx44\ \text{klaszter}$}\end{center}
\jegyzet{számérték}{}

\lepes{b) Egyszerre beszélők}
Egy klaszter mind a 124 vivőt felhasználja, mindegyiken 8 időrés, egy időrés egy beszélgetés.

\begin{center}\adjustbox{max width=\linewidth}{$\displaystyle 124\cdot8=992\ \text{beszélgetés klaszterenként}$}\end{center}
\jegyzet{számérték}{}

\begin{center}\adjustbox{max width=\linewidth}{$\displaystyle 43{,}75\cdot992=43\,400$}\end{center}
\jegyzet{számérték}{kb. 43 400 egyszerre a városban}

\lepes{c) Egy cella, egy szolgáltató}
A szolgáltató 40 vivője 12 cella között oszlik meg.

\begin{center}\adjustbox{max width=\linewidth}{$\displaystyle \frac{40}{12}=3{,}33\ \Rightarrow\ 3\ \text{vagy}\ 4\ \text{vivő cellánként}$}\end{center}
\jegyzet{számérték}{}

\begin{center}\adjustbox{max width=\linewidth}{$\displaystyle 3\cdot8=24\ \text{és}\ 4\cdot8=32\ \text{fő}$}\end{center}
\jegyzet{számérték}{}
\begin{eredmeny}

\begin{center}\adjustbox{max width=\linewidth}{$\displaystyle \text{a) }\approx44\ \text{klaszter}\qquad\text{b) }\approx43\,400\ \text{beszélgetés}\qquad\text{c) }24\ldots32\ \text{fő}$}\end{center}
\end{eredmeny}
\begin{tipp}[Ellenőrzés]
\begin{itemize}
\item A gyűjtemény szerint több hely: Half Rate, AMR, vagy az 1800 MHz-es sáv (75 MHz, 374 vivő). A szolgáltatónkénti vivőszám régi adat \texttt{[ELLENŐRIZENDŐ]}.
\item Valójában néhány időrés jelzésre megy el, ezért ez felső becslés.
\end{itemize}
\end{tipp}


% ------------------------------------------------------------------
\section{\texorpdfstring{9.2. feladat: Roaming: hívás, bontás, SMS}{9.2. feladat: Roaming: hívás, bontás, SMS}}\label{o9_02}
\setcounter{lepes}{0}
\begin{feladat}{9.2. feladat}
MS-A az $\alpha$, MS-B a $\beta$ előfizetője, MS-A $\beta$-ban, MS-B $\alpha$-ban barangol. MS-A felhívja MS-B-t. a) Melyik HLR? b) Melyik elem ellenőrzi a hívó készülék tiltását? MS-B bont: c) ki értesül először? d) kik értesülnek? MS-A SMS-t küld MS-B-nek: e) melyik SMSC kapja? f) melyik kézbesíti?

{\footnotesize\sffamily Forrás: 5-ösöm 27. oldal $\cdot$ 5o\_27\_1. Címkék: vizsga, javítás}
\end{feladat}
\begin{keresunk}
\begin{itemize}
\item a) a lekérdezett HLR
\item b) az EIR
\item c), d) a bontásról értesülők
\item e), f) az SMSC
\end{itemize}
\end{keresunk}
\begin{jeloles}
\begin{jeltabla}
$\beta$ & a másik szolgáltató ($\beta$) hálózata, alsó indexben: a $\beta$ hálózat eleme & \\
$\alpha$ & az egyik szolgáltató ($\alpha$) hálózata, alsó indexben: az $\alpha$ hálózat eleme & \\
\end{jeltabla}
\par\smallskip\textbf{Adott még:}
\begin{itemize}
\item MS-A: $\alpha$-előfizető, a $\beta$ hálózatban van
\item MS-B: $\beta$-előfizető, az $\alpha$ hálózatban van
\end{itemize}
\end{jeloles}

\lepes{A hívás útja}
MS-A a $\beta$ hálózatban indítja a hívást: BTS$_\beta$ $\to$ BSC$_\beta$ $\to$ MSC$_\beta$. A hívott szám $\beta$-beli, ezért a hívás a GMSC$_\beta$-hoz fut, az lekérdezi a HLR$_\beta$-t, amely tudja, hogy MS-B a VLR$_\alpha$ területén van. A hívás innen MSC$_\alpha$ $\to$ BSC$_\alpha$ $\to$ BTS$_\alpha$ $\to$ MS-B úton halad.

\lepes{a) HLR}
1. szabály: a \textbf{hívott} (MS-B) honi HLR-je.

\begin{center}\adjustbox{max width=\linewidth}{$\displaystyle \text{HLR}_\beta$}\end{center}
\jegyzet{definíció}{MS-B $\beta$-előfizető}

\lepes{b) A készülék ellenőrzése}
3. szabály: a hívó készülékét az a hálózat ellenőrzi, ahol a készülék éppen van. MS-A a $\beta$-ban van.

\begin{center}\adjustbox{max width=\linewidth}{$\displaystyle \text{MSC}_\beta\ \text{az EIR}_\beta\text{-nél}$}\end{center}
\jegyzet{definíció}{}

\lepes{c) Ki értesül először a bontásról?}
5. szabály: MS-B bont, és ő a BTS$_\alpha$-val beszél közvetlenül.

\begin{center}\adjustbox{max width=\linewidth}{$\displaystyle \text{BTS}_\alpha$}\end{center}
\jegyzet{definíció}{}

\lepes{d) Kik értesülnek?}
A bontás a teljes összeköttetésen végigfut, minden elem felszabadítja az erőforrásait: a hívás útjának minden eleme.

\begin{center}\adjustbox{max width=\linewidth}{$\displaystyle \text{BTS}_\alpha,\ \text{BSC}_\alpha,\ \text{MSC}_\alpha,\ \text{GMSC}_\beta,\ \text{MSC}_\beta,\ \text{BSC}_\beta,\ \text{BTS}_\beta$}\end{center}
\jegyzet{definíció}{és végül MS-A}

\lepes{e), f) SMS}
2. szabály: a \textbf{küldő} (MS-A) honi SMSC-je fogadja és kézbesíti, akárhol van is MS-A.

\begin{center}\adjustbox{max width=\linewidth}{$\displaystyle \text{SMSC}_\alpha$}\end{center}
\jegyzet{definíció}{mindkét kérdésre}
\begin{tipp}[A gyűjtemény megoldása]
A gyűjtemény b)-re EIR$_\alpha$-t ír. Ugyanezen az oldalon egy másik változatban (MS-B SMS-e MS-A-nak) SMSC$_\alpha$-t.
\end{tipp}
\begin{figyelem}[Javítás]
b) MS-A a $\beta$ hálózatban van, a kiszolgáló hálózat ellenőriz: EIR$_\beta$ \texttt{[ELLENŐRIZENDŐ]}. A másik változatban a küldő MS-B, az ő honi SMSC-je az SMSC$_\beta$.
\end{figyelem}
\begin{eredmeny}

\begin{center}\adjustbox{max width=\linewidth}{$\displaystyle \text{a) HLR}_\beta\quad\text{b) EIR}_\beta\quad\text{c) BTS}_\alpha\quad\text{d) a hívás útjának minden eleme}\quad\text{e), f) SMSC}_\alpha$}\end{center}
\end{eredmeny}


% ------------------------------------------------------------------
\section{\texorpdfstring{9.3. feladat: Két $\alpha$-előfizető, az egyik $\beta$-ban}{9.3. feladat: Két α-előfizető, az egyik β-ban}}\label{o9_03}
\setcounter{lepes}{0}
\begin{feladat}{9.3. feladat}
MS-A1 és MS-A2 is $\alpha$-előfizető, MS-A2 $\beta$-ban barangol. a) MS-A1 hívja MS-A2-t: melyik VLR-t kérdezi a HLR$_\alpha$? b) Melyik MSC értesül? c) Ki tud MS-A2 beszélgetés közbeni mozgásáról? d) Bontás után MS-A2 SMS-t ír a készülékét a tengerbe dobó MS-A1-nek: ki tárolja?

{\footnotesize\sffamily Forrás: 5-ösöm 28. oldal $\cdot$ 5o\_28\_1. Címkék: vizsga}
\end{feladat}
\begin{keresunk}
\begin{itemize}
\item a) a VLR
\item b) az MSC
\item c) a mozgásról tudó elemek
\item d) az SMS-t tároló elem
\end{itemize}
\end{keresunk}
\begin{jeloles}
\begin{jeltabla}
$\alpha$ & az egyik szolgáltató ($\alpha$) hálózata, alsó indexben: az $\alpha$ hálózat eleme & \\
$\beta$ & a másik szolgáltató ($\beta$) hálózata, alsó indexben: a $\beta$ hálózat eleme & \\
\end{jeltabla}
\par\smallskip\textbf{Adott még:}
\begin{itemize}
\item MS-A1: $\alpha$-előfizető, otthon ($\alpha$)
\item MS-A2: $\alpha$-előfizető, a $\beta$ hálózatban
\end{itemize}
\end{jeloles}

\lepes{a) VLR}
A hívott (MS-A2) $\alpha$-előfizető, ezért a HLR$_\alpha$-t kérdezik (1. szabály). A HLR$_\alpha$ tudja, hogy MS-A2 a $\beta$ hálózat VLR-jénél van bejelentkezve, ettől kér barangolási számot.

\begin{center}\adjustbox{max width=\linewidth}{$\displaystyle \text{VLR}_\beta$}\end{center}
\jegyzet{definíció}{}

\lepes{b) MSC}
A hívást MS-A1 indítja $\alpha$-ban, így elsőként a saját kiszolgáló MSC-je értesül. A hívás ezután az MSC$_\beta$-hoz jut, amely MS-A2-t kiszolgálja.

\begin{center}\adjustbox{max width=\linewidth}{$\displaystyle \text{MSC}_\alpha,\ \text{majd MSC}_\beta$}\end{center}
\jegyzet{definíció}{a kérdés kétértelmű}

\lepes{c) Mozgás beszélgetés közben}
4. szabály: hívásátadás, csak a kiszolgáló hálózat rádiós és kapcsoló elemei vesznek részt benne.

\begin{center}\adjustbox{max width=\linewidth}{$\displaystyle \text{BTS}_\beta,\ \text{BSC}_\beta,\ \text{MSC}_\beta$}\end{center}
\jegyzet{definíció}{}

\lepes{d) SMS tárolása}
2. szabály: a küldő (MS-A2) honi SMSC-je fogadja. Mivel MS-A1 nem érhető el, ott várakozik, amíg kézbesíthető nem lesz.

\begin{center}\adjustbox{max width=\linewidth}{$\displaystyle \text{SMSC}_\alpha$}\end{center}
\jegyzet{definíció}{}
\begin{eredmeny}

\begin{center}\adjustbox{max width=\linewidth}{$\displaystyle \text{a) VLR}_\beta\quad\text{b) MSC}_\alpha\ (\text{a hívóé}),\ \text{majd MSC}_\beta\quad\text{c) BTS}_\beta,\ \text{BSC}_\beta,\ \text{MSC}_\beta\quad\text{d) SMSC}_\alpha$}\end{center}
\end{eredmeny}
\begin{tipp}[Ellenőrzés]
\begin{itemize}
\item c) A gyűjtemény HLR$_\alpha$-t és VLR$_\beta$-t is ír, hívás közben azonban nincs helyzetfrissítés. A b) kérdés kétértelmű \texttt{[ELLENŐRIZENDŐ]}.
\end{itemize}
\end{tipp}


% ------------------------------------------------------------------
\section{\texorpdfstring{9.4. feladat: AuC, FDM és TDM, VAD, a betűszavak}{9.4. feladat: AuC, FDM és TDM, VAD, a betűszavak}}\label{o9_04}
\setcounter{lepes}{0}
\begin{feladat}{9.4. feladat}
Mi az AuC szerepe? Mit jelent, hogy a GSM FDM-et és TDM-et is használ? Mi a VAD előnye GSM-ben és VoIP-ban, és hogyan tartják fenn a kapcsolat érzetét? Ismertesse: HLR, VLR, EIR, MSC, BSC, BTS, FNR.

{\footnotesize\sffamily Forrás: 5-ösöm 21., 27\ldots{}28. oldal $\cdot$ 5o\_21\_2, 5o\_27\_1, 5o\_28\_1. Címkék: beugró}
\end{feladat}
\begin{keresunk}
\begin{itemize}
\item rövid, indokolt válaszok
\end{itemize}
\end{keresunk}
\begin{jeloles}
\par\smallskip\textbf{Adott még:}
\begin{itemize}
\item fogalmi kérdések a GSM hálózat elemeiről és működéséről
\end{itemize}
\end{jeloles}

\lepes{AuC}
A SIM titkos kulcsát (Ki) tárolja. Hitelesítéskor véletlen számot (RAND) küld a készüléknek, és a Ki-ből az A3 algoritmussal kiszámolja a várt választ (SRES). Ha a SIM ugyanezt adja, a SIM valódi. Az A8 algoritmussal ugyanebből a rejtjelező kulcs (Kc) is előáll, ezzel titkosítják a rádiós szakaszt (A5).

\lepes{FDM és TDM}
FDM: a sáv 200 kHz-es vivőkre van osztva. TDM: minden vivőn egy keret 8 időrésből áll, mindegyik időrés egy-egy beszélgetésé. Egy beszélgetést tehát egy vivő és egy időrés jelöl ki.

\lepes{VAD}
A hangaktivitás-érzékelő (VAD) felismeri a beszédszünetet, ilyenkor nem adnak. GSM-ben ez kisebb akkumulátor-fogyasztást és kisebb interferenciát jelent, VoIP-ban kevesebb csomagot, kisebb sávot. Hogy a vevő ne higgye a vonalat megszakadtnak, a vevőoldal \textbf{komfortzajt} generál (a szünet elején átküldött zajjellemzők alapján).

\lepes{A betűszavak}
HLR: honi előfizetői adatbázis. VLR: a területen tartózkodók ideiglenes adatai. EIR: IMEI-listák, fehér (engedélyezett), fekete (tiltott), szürke (megfigyelt). MSC: kapcsolóközpont, hívásirányítás. BSC: bázisállomás-vezérlő. BTS: rádiós adó-vevő. FNR: számhordozhatósági adatbázis (melyik szolgáltatóhoz került át egy szám).
\begin{eredmeny}

\begin{center}\adjustbox{max width=\linewidth}{$\displaystyle \text{AuC: hitelesítés (Ki, RAND, SRES, Kc)}\quad\text{200 kHz-es vivők, 8 időrés}\quad\text{VAD: kevesebb adás, komfortzaj}$}\end{center}
\end{eredmeny}
